\section{Autoencoder-Based Representation Learning in Alzheimer's Disease}

reference papers:
\begin{itemize}
    \item “Deep Learning-Based Feature Representation for AD/MCI Classification” — Heung-Il Suk \& Dinggang Shen (2013). Deep learning for AD/MCI representation learning, extraction of disease-relevant features from neuroimaging/biomarker data, use of learned representations for diagnostic classification.
    \item “Reducing Variations in Multi-Center Alzheimer's Disease Classification with Convolutional Adversarial Autoencoder” — Bernard M. Cobbinah et al. (2022). Convolutional adversarial AE for Alzheimer's classification, learning disease-relevant representations while reducing differences caused by multi-center data, downstream diagnostic classification.
    \item “Automated Prediction System for Alzheimer Detection Based on Deep Residual Autoencoder and Support Vector Machine” — M. Menagadevi, S. Mangai, Nirmala Madian \& D. Thiyagarajan (2023). Deep residual AE for Alzheimer's feature extraction, dimensionality reduction followed by SVM classification, use of AE representations for automated AD detection.
    \item “Latent Diffusion Autoencoders: Toward Efficient and Meaningful Unsupervised Representation Learning in Medical Imaging – A Case Study on Alzheimer's Disease” — Gabriele Lozupone, Alessandro Bria, Francesco Fontanella, Frederick J. A. Meijer, Claudio De Stefano \& Henkjan Huisman (2026). Latent autoencoder representation learning for medical imaging, Alzheimer's disease case study, focus on learning compact and meaningful unsupervised representations.
    \item “Using Deep Autoencoders to Identify Abnormal Brain Structural Patterns in Neuropsychiatric Disorders: A Large-Scale Multi-Sample Study” — Walter H. L. Pinaya, Andrea Mechelli \& João R. Sato (2018). Deep AE applied to structural brain imaging, learning normative representations of healthy brain structure, use of reconstruction differences to identify abnormal patterns in clinical populations. MAYBE as it is not just AD but other disorders too.
\end{itemize}
